\documentclass[11pt]{article}
\usepackage[margin=1.05in]{geometry}
\usepackage{amsmath,amssymb,amsthm,mathtools}
\usepackage{microtype}
\usepackage{booktabs}
\usepackage{enumitem}
\usepackage[hidelinks]{hyperref}
\usepackage[nameinlink,noabbrev]{cleveref}

\newtheorem{theorem}{Theorem}[section]
\newtheorem{proposition}[theorem]{Proposition}
\newtheorem{lemma}[theorem]{Lemma}
\newtheorem{corollary}[theorem]{Corollary}
\theoremstyle{definition}
\newtheorem{definition}[theorem]{Definition}
\theoremstyle{remark}
\newtheorem{remark}[theorem]{Remark}

\newcommand{\SSer}{\mathfrak S}
\newcommand{\Mob}{\mathsf M}
\newcommand{\ZE}{\mathsf Z_E}
\newcommand{\evE}{\operatorname{ev}_E}
\newcommand{\1}{\mathbf 1}
\newcommand{\e}{\mathrm e}
\newcommand{\ph}{\varphi}
\newcommand{\Om}{\Omega}
\newcommand{\Dcal}{\mathcal D}
\newcommand{\Fcal}{\mathcal F}
\newcommand{\Kcal}{\mathcal K}
\newcommand{\Lcal}{\mathcal L}
\newcommand{\Zmod}{\mathbb Z}
\newcommand{\eps}{\varepsilon}

\title{An Exact Contraction Dictionary and a Fixed-Modulus Anchored Transfer Theorem for Singular Series}
\author{Charlie Dorsey}
\date{September 2026}

\begin{document}
\maketitle

\begin{abstract}
Fix a squarefree modulus $q$ and compatible reduced endpoint classes $a,b\pmod q$.  We organize the endpoint-anchored and free connected singular-series coordinates by a common Boolean contraction calculus, identify the completed endpoint coordinate exactly with the free-label connected coordinate, and isolate the anchored-to-free defect as the top two-anchor connected contraction residual.  We then prove, pointwise along compatible heights,
\[
A_{1,q}(a,b;h)-B_{1,q}(a,b;h)=o_q(\log h).
\]
The fixed-$q$ proof isolates the finite-period changes in the pointwise $q=1$ argument: the low-prime and admissibility data form a fixed periodic decoration, the high-prime Boolean Fourier coefficients have unchanged absolute value, the exterior completion and reciprocal estimates incur only fixed-conductor constants, and the one-dimensional logarithmic main terms cancel before the nonprincipal periodic corrections.  We also record a finite-cutoff arithmetic diagonal-contraction identity for the active high-prime coefficient system.  No rank-uniform conclusion is asserted.
\end{abstract}

\section{Introduction and scope}\label{sec:intro}

For a finite set $H\subset\Zmod$ and a squarefree integer $q\ge1$, define the $q$-restricted singular series
\begin{equation}\label{eq:S_q}
\SSer_q(H)
=\prod_{p\nmid q}
\left(1-\frac{\nu_p(H)}p\right)
\left(1-\frac1p\right)^{-|H|},
\end{equation}
where $\nu_p(H)$ is the number of residue classes occupied by $H$ modulo $p$.  As usual, the product is zero when the tuple is inadmissible at an active prime.  Let
\begin{equation}\label{eq:S_q0}
\SSer_{q,0}(H)
=\sum_{J\subseteq H}(-1)^{|H|-|J|}\SSer_q(J)
\end{equation}
be its Boolean M\"obius transform.

The purpose of this note is narrow.  First, we give an exact finite dictionary between endpoint completion and free-label connected coordinates.  Second, we prove the fixed-squarefree-$q$ pointwise level-one anchored-to-free transfer theorem.  Third, we record a finite-cutoff arithmetic diagonal-contraction identity for the active high-prime coefficient system.  We do not address an all-orders tail, a growing-rank theorem, or any resummation problem.

At level one, the two quantities compared by the theorem have simple roles.  The sum $A_{1,q}$ measures the connected contribution of one admissible interior point relative to the two fixed endpoints, while $B_{1,q}$ is the sum of the two one-anchor connected boundary contributions.  Their exact formulas appear in \eqref{eq:A1q}--\eqref{eq:B1q}; the theorem proves that their difference is $o_q(\log h)$ for every compatible height progression.

The proof has four stages.  The Boolean contraction identity in \cref{sec:dictionary} first identifies the correct connected coordinate exactly.  The fixed-period reduction in \cref{sec:fixed-period} then isolates the only changes caused by the modulus $q$: a bounded low-period decoration and the deletion of finitely many Euler factors.  The analytic work is concentrated in two places: the Boolean small box, where the companion $q=1$ argument imports Shiu's short-interval estimate, and the exterior, where weighted completion, reciprocity, and Karatsuba's trigonometric-sum estimate are used.  The remaining one-dimensional terms are matched by a fixed-conductor P\'etermann estimate.  Each imported result is stated below with the specific estimate and hypotheses used; the complete $q=1$ derivations are carried by the companion manuscript \cite{DorseyPointwise2026} rather than repeated here.

Fixed physical anchors and the conditional fixed-gap Poisson factor have established antecedents in the consecutive-prime literature \cite{GoldstonLedoan2011,LOS}; congruence-restricted singular-series sums likewise have substantial prior art \cite{Kuperberg2025}.  So do Boolean M\"obius inversion and generic connected/cumulant language.  No novelty claim for those ingredients is made here.  The priority posture of the present theorem is stated separately in \cref{sec:priority}; mathematical statements below do not depend on that posture.

\section{The connected contraction dictionary}\label{sec:dictionary}

Fix the endpoint set
\[
E=\{0,h\}.
\]
Let $T$ be any finite set of free labels disjoint from $E$.  For $A\subseteq E$ and $U\subseteq T$, put
\[
f_q(A,U)=\SSer_q(A\cup U).
\]
Define partial Boolean M\"obius operators
\begin{align}
(\Mob_E f)(A,U)
&=\sum_{B\subseteq A}(-1)^{|A|-|B|}f(B,U),\label{eq:ME}\\
(\Mob_T f)(A,U)
&=\sum_{V\subseteq U}(-1)^{|U|-|V|}f(A,V).\label{eq:MT}
\end{align}
They commute, and
\[
(\Mob_E\Mob_Tf_q)(A,U)=\SSer_{q,0}(A\cup U).
\]
Let endpoint-fiber contraction be
\[
(\ZE g)(U)=\sum_{A\subseteq E}g(A,U),
\]
and let $\evE$ denote evaluation at the full endpoint set $A=E$.

\begin{proposition}[Endpoint contraction identity]\label{prop:endpoint-contraction}
For every $q$, every endpoint pair $E$, and every finite free set $T$,
\begin{equation}\label{eq:commuting-identity}
\boxed{\ZE\Mob_E\Mob_T=\Mob_T\evE.}
\end{equation}
Equivalently,
\begin{equation}\label{eq:F_E}
\Fcal_{E,q}(T)
:=\sum_{A\subseteq E}\SSer_{q,0}(A\cup T)
=\Mob_T[\SSer_q(E\cup\cdot)](T).
\end{equation}
\end{proposition}

\begin{proof}
For fixed $U\subseteq T$,
\[
\sum_{A\subseteq E}\sum_{B\subseteq A}
(-1)^{|A|-|B|}f_q(B,U)
=f_q(E,U),
\]
which is the Boolean zeta--M\"obius identity on the endpoint fiber.  Since $\Mob_E$ and $\Mob_T$ commute, \eqref{eq:commuting-identity} follows.
\end{proof}

Assume $\SSer_q(E)>0$ and define
\begin{equation}\label{eq:Rq}
R_q(U)=\frac{\SSer_q(E\cup U)}{\SSer_q(E)},
\qquad
R_q^\circ(T)=\sum_{V\subseteq T}(-1)^{|T|-|V|}R_q(V).
\end{equation}
Then \cref{prop:endpoint-contraction} gives the pointwise identity
\begin{equation}\label{eq:F-Rcirc}
\boxed{\Fcal_{E,q}(T)=\SSer_q(E)R_q^\circ(T).}
\end{equation}
If $I$ is a finite interior set and
\[
C_{k,q}^{\mathrm{comp}}
=\sum_{\substack{T\subseteq I\\|T|=k}}\Fcal_{E,q}(T),
\qquad
A_{k,q}^{\circ}
=\sum_{\substack{T\subseteq I\\|T|=k}}R_q^\circ(T),
\]
then
\begin{equation}\label{eq:C-A-dictionary}
\boxed{C_{k,q}^{\mathrm{comp}}=\SSer_q(E)A_{k,q}^{\circ}.}
\end{equation}
At $q=1$ this is the exact normalization dictionary between the completed endpoint coordinate and the anchored free-label connected coordinate.

\subsection{Anchor subtraction}

Write
\[
c_q(A,T)=\SSer_{q,0}(A\cup T),
\]
and define the lower-anchor sum and the top-anchor sector by
\begin{align}
b_q(T)&=c_q(\varnothing,T)+c_q(\{0\},T)+c_q(\{h\},T),\label{eq:bq}\\
g_q(T)&=c_q(E,T)=\SSer_{q,0}(E\cup T).\label{eq:gq}
\end{align}
Then $\Fcal_{E,q}(T)=b_q(T)+g_q(T)$.  Since singleton singular series are one,
\[
\SSer_{q,0}(E)=\SSer_q(E)-1.
\]
Combining this with \eqref{eq:F-Rcirc} gives the exact residual identity
\begin{equation}\label{eq:anchor-residual}
\boxed{
R_q^\circ(T)-b_q(T)
=\frac{g_q(T)-\SSer_{q,0}(E)b_q(T)}{\SSer_q(E)}.
}
\end{equation}
Thus the anchored-to-free transfer defect is the normalized failure of the top two-anchor connected sector to contract onto the lower anchor sectors.

\subsection{The same transform at one prime}

At one active prime, write
\[
L_p(H)=\mathbf E_z\prod_{x\in H}\phi_z(x),
\]
where $\phi_z(x)$ is the normalized local avoidance factor.  If a coordinate value is repeated, $H$ is read as a labeled multiset and the product is over coordinate occurrences.  Boolean connectedization gives
\begin{equation}\label{eq:local-connected}
\sum_{J\subseteq H}(-1)^{|H|-|J|}L_p(J)
=\mathbf E_z\prod_{x\in H}(\phi_z(x)-1).
\end{equation}
This is the local-prime image of the same connected transform: every free Fourier frequency is forced away from zero after connectedization.

\section{Fixed squarefree modulus and the level-one theorem}\label{sec:fixed-q-setup}

Fix squarefree $q\ge1$ and endpoint classes
\[
a,b\in(\Zmod/q\Zmod)^\times,
\qquad h\equiv b-a\pmod q.
\]
Assume the endpoint pair is $q$-restricted admissible,
\[
\SSer_q(E)>0,
\qquad E=\{0,h\}.
\]
Every legal height is even: if $2\nmid q$, admissibility at the active prime $2$ forces $h$ even; if $2\mid q$, the unit classes $a,b$ are both odd and hence their difference is even modulo $2$, so $h\equiv b-a\equiv0\pmod2$.

Define the admissible interior and its exact size by
\begin{equation}\label{eq:IqNq}
I_q(a;h)=\{0<t<h:(a+t,q)=1\},
\qquad N_q(a;h)=|I_q(a;h)|.
\end{equation}
The level-one connected and lower-anchor sums are
\begin{align}
A_{1,q}(a,b;h)
&=\sum_{t\in I_q(a;h)}R_q^\circ(\{t\})\notag\\
&=\sum_{t\in I_q(a;h)}
\left(
\frac{\SSer_q(\{0,t,h\})}{\SSer_q(\{0,h\})}-1
\right),\label{eq:A1q}\\
B_{1,q}(a,b;h)
&=\sum_{t\in I_q(a;h)}b_q(\{t\})\notag\\
&=\sum_{t\in I_q(a;h)}
\left[\SSer_{q,0}(\{t\})
+\SSer_{q,0}(\{0,t\})
+\SSer_{q,0}(\{h,t\})\right].\label{eq:B1q}
\end{align}
Since $\SSer_{q,0}(\{t\})=0$,
\begin{equation}\label{eq:B1q-simple}
B_{1,q}
=\sum_{t\in I_q(a;h)}(\SSer_q(\{0,t\})-1)
+\sum_{t\in I_q(a;h)}(\SSer_q(\{h,t\})-1).
\end{equation}
By \eqref{eq:anchor-residual},
\begin{equation}\label{eq:AminusB-connected}
A_{1,q}-B_{1,q}
=\frac1{\SSer_q(E)}
\sum_{t\in I_q(a;h)}
\left[g_q(\{t\})-\SSer_{q,0}(E)b_q(\{t\})\right].
\end{equation}

\begin{theorem}[Fixed-squarefree-$q$ pointwise level-one transfer]\label{thm:fixed-q}
For every fixed squarefree $q$, uniformly over the finitely many compatible admissible endpoint residue pairs $(a,b)$,
\begin{equation}\label{eq:main-theorem}
\boxed{A_{1,q}(a,b;h)-B_{1,q}(a,b;h)=o_q(\log h)}
\end{equation}
as $h\to\infty$ through positive compatible heights.
\end{theorem}

There is no density-one restriction in $h$.

\subsection{Exact density normalization}

The admissible count is exactly periodic.

\begin{lemma}[Exact periodic count]\label{lem:Nq}
For each compatible pair $(a,b)$,
\begin{equation}\label{eq:Nq-periodic}
\boxed{N_q(a;h)=\frac{\ph(q)}q h+\eps_q(a,b),}
\end{equation}
where $\eps_q(a,b)$ depends only on the fixed compatible residue data, not on the positive representative of the gap class.
\end{lemma}

\begin{proof}
If $h$ is increased by $q$ while preserving its residue class, the interval gains one complete period of $q$ translates and hence exactly $\ph(q)$ admissible interior positions.  Therefore $N_q(a;h)-\ph(q)h/q$ is constant on each compatible height progression.
\end{proof}

For the completed Lemke Oliver--Soundararajan normalization, the fixed-modulus density parameter is
\begin{equation}\label{eq:deltaq}
\delta_q(x)=\frac{q}{\ph(q)\log x}.
\end{equation}
Consequently
\[
\delta_q(x)\frac{\ph(q)}q=\frac1{\log x}.
\]
The $q/\ph(q)$ factor belongs to the ambient admissible-point density; it does not appear as an additional multiplier inside the transfer kernel \eqref{eq:main-theorem}.  We use $\Delta\mu_0$ only as notation for the completed level-one coefficient per admissible interior site,
\[
\Delta\mu_0:=\frac{C_{1,q}^{\mathrm{comp}}}{N_q(a;h)}.
\]
Thus \eqref{eq:F-Rcirc} gives at level one
\begin{equation}\label{eq:LOS-level-one}
C_{1,q}^{\mathrm{comp}}=\SSer_q(E)A_{1,q},
\qquad
\Delta\mu_0=\frac{\SSer_q(E)}{N_q(a;h)}A_{1,q}.
\end{equation}

\section{Exact fixed-period reduction}\label{sec:fixed-period}

The pointwise $q=1$ proof is carried by the companion estimate
\begin{equation}\label{eq:q1-input}
\sum_{0<t<h}\left(
\frac{\SSer(\{0,t,h\})}{\SSer(\{0,h\})}-1
\right)
-2\sum_{0<t<h}(\SSer(\{0,t\})-1)
=o(\log h)
\end{equation}
for even $h$.  We now isolate every change caused by fixed $q$ before invoking its analytic bounds.

Put
\[
\Lambda_q=\operatorname{lcm}(6,q),
\qquad
\operatorname{rad}_{\ge5}(n):=\prod_{\substack{p\mid n\\p\ge5}}p.
\]
We write $\e(t):=\exp(2\pi i t)$ and, for an integer modulus $M\ge1$,
$\e_M(x):=\e(x/M)$.  For the anchored ratio, collect all admissibility and active low-prime data into the bounded $\Lambda_q$-periodic weight
\begin{equation}\label{eq:lambdaq}
\lambda_{q,a,h}(y)
=\1_{(a+y,q)=1}
\prod_{\substack{p\in\{2,3\}\\p\nmid q}}
\frac{L_p(\{0,y,h\})}{L_p(\{0,h\})}.
\end{equation}
All remaining divisor variables are supported on primes $p\ge5$ with $p\nmid q$.

For $p\ge5$ set
\begin{equation}\label{eq:c0c1}
c_0(p)=\frac{p(p-3)}{(p-1)(p-2)},
\qquad
c_1(p)=\frac{p(p-2)}{(p-1)^2}.
\end{equation}
Define squarefree coefficient systems
\begin{align}
\omega_q(e)
&=\mu^2(e)\1_{(e,6q)=1}\prod_{p\mid e}\frac1{p-3},\label{eq:omegaq}\\
\eta_q(e)
&=\mu^2(e)\1_{(e,6q)=1}\prod_{p\mid e}\frac1{p-2},\label{eq:etaq}\\
v_{q,h}(e)
&=\mu^2(e)\1_{(e,6q)=1}
\prod_{\substack{p\mid e\\p\nmid h}}\frac1{p-3}
\prod_{\substack{p\mid e\\p\mid h}}\frac1{p-2},\label{eq:vqh}\\
C_q^*(h)
&=\prod_{\substack{p\ge5\\p\nmid qh}}c_0(p)
\prod_{\substack{p\ge5\\p\nmid q,\ p\mid h}}c_1(p).\label{eq:Cqstar}
\end{align}
For coprime high divisors, put
\begin{align}
N_{q,a}(h;e_1,e_2)
&=\sum_{\substack{0<y<h\\e_1\mid y,\ e_2\mid h-y}}
\lambda_{q,a,h}(y),\label{eq:Nqea}\\
\Delta_{q,a}(h;e_1,e_2)
&=N_{q,a}(h;e_1,e_2)-\frac{N_q(a;h)}{e_1e_2}.\label{eq:Deltaqea}
\end{align}

\begin{proposition}[Exact divisor expansion]\label{prop:divisor-expansion}
The anchored level-one sum admits the exact identity
\begin{equation}\label{eq:Tq-expansion}
\boxed{
A_{1,q}(a,b;h)
=C_q^*(h)
\sum_{\substack{e_1,e_2\ge1\\(e_1,e_2)=1,\ (e_2,h)=1}}
v_{q,h}(e_1)\omega_q(e_2)
\Delta_{q,a}(h;e_1,e_2).
}
\end{equation}
Every divisor on the right is supported on primes $p\ge5$ with $p\nmid q$.
\end{proposition}

\begin{proof}
Expand the active high-prime local ratio prime by prime.  Keep the fixed low-prime and admissibility factor \eqref{eq:lambdaq} unexpanded.  At each active high prime, the formal baseline together with its permitted selector or selectors sums to one.  Primes dividing $q$ occur in neither local singular-series factor nor the high divisor systems.  Thus insertion of the exact baseline $N_q(a;h)/(e_1e_2)$ reconstructs the literal subtraction $-N_q(a;h)$ in \eqref{eq:A1q}; no asymptotic density replacement is used.
\end{proof}

\subsection{Rowwise normalization}\label{sec:row-contract}
The two-dimensional estimates are applied row by row, and the actual row
height must remain visible.  We write $H$ for the parity-reduced kernel
height: $H=h/2$ on an ordinary row and $H=h/(2d)$ on a row carrying a
deterministic high divisor $d\mid h$.  Every progression root and oscillatory
phase in that row uses this $H$.  Let
$R\mid\operatorname{rad}_{\ge5}(h)$ be the squarefree product of additional
forbidden high height primes inherited from the row reduction, kept separate
from the fixed modulus $q$; active fixed-product moduli satisfy $(m,HqR)=1$.

For one exact low-period row, let $\lambda$ be its bounded periodic weight,
with period dividing $\Lambda_q$, and define
\begin{align}
N_\lambda(H)&=\sum_{0<y<H}\lambda(y),\label{eq:Nlambda}\\
D_\lambda(e_1,e_2;H)
&=\sum_{\substack{0<y<H\\e_1\mid y,\ e_2\mid H-y}}\lambda(y)
 -\frac{N_\lambda(H)}{e_1e_2}.\label{eq:Dlambda}
\end{align}

\begin{lemma}[Exact fixed-$q$ row reduction]\label{lem:fixed-q-row-reduction}
For a row of \eqref{eq:Tq-expansion}, write
\[
d=(e_1,h),\qquad r=e_1/d.
\]
Thus $d$ is either $1$ or a squarefree divisor of
$\operatorname{rad}_{\ge5}(h)$ supported away from $q$, and put
\[
h_d=\frac hd=2H_d,
\qquad
 a_d\equiv d^{-1}a\pmod q,
\]
and define the reduced low-period weight
\begin{equation}\label{eq:reduced-low-weight-q}
\lambda^{[d]}_{q,a,h}(z):=\lambda_{q,a_d,h_d}(2z).
\end{equation}
For every such row one has the exact identity
\begin{equation}\label{eq:fixed-q-row-reduction}
\boxed{
\Delta_{q,a}(h;dr,e_2)
=D_{\lambda^{[d]}_{q,a,h}}(r,e_2;H_d)
+\frac{\delta^{\rm per}_{q,a,h}(d)}{r e_2},
}
\end{equation}
where
\begin{equation}\label{eq:periodic-density-delta-q}
\delta^{\rm per}_{q,a,h}(d)
:=N_{\lambda^{[d]}_{q,a,h}}(H_d)-\frac{N_q(a;h)}d
=O_q(1)
\end{equation}
uniformly in the permitted $d$ and compatible heights.  Consequently the
periodic-density term in \eqref{eq:fixed-q-row-reduction} contributes
$O_q(1)$ on the ordinary two-dimensional family and
$O_q(\log\log h)$ after the outer decorated-row sum.
\end{lemma}

\begin{proof}
Write $y=dz$.  Since every prime of $d$ is at least $5$ and $d$ is a unit
modulo $6q$, multiplication by $d$ preserves the active $2,3$ local ratio and
sends the admissibility condition to $(a_d+z,q)=1$.  Hence, literally,
\[
\lambda_{q,a,h}(dz)=\lambda_{q,a_d,h_d}(z).
\]
The legal height $h_d$ is even.  If $2\nmid q$, its active $p=2$ factor kills
odd $z$; if $2\mid q$, the unit residue $a_d$ is odd and the admissibility
indicator does the same.  Thus $z=2x$ and, because $r,e_2$ are odd,
\[
\sum_{\substack{0<y<h\\dr\mid y,\ e_2\mid h-y}}
\lambda_{q,a,h}(y)
=
\sum_{\substack{0<x<H_d\\r\mid x,\ e_2\mid H_d-x}}
\lambda^{[d]}_{q,a,h}(x).
\]
Subtracting $N_q(a;h)/(dr e_2)$ on the left and
$N_{\lambda^{[d]}_{q,a,h}}(H_d)/(r e_2)$ on the right proves
\eqref{eq:fixed-q-row-reduction}.

For the bound, one complete $\Lambda_q$-period of
$z\mapsto\lambda_{q,a_d,h_d}(2z)$ has mean $2\ph(q)/q$: the admissibility
indicator contributes $\ph(q)/q$, and each active low-prime local ratio has
complete-period mean one before the parity compression.  Hence
\[
N_{\lambda^{[d]}_{q,a,h}}(H_d)
=\frac{\ph(q)}q h_d+O_q(1).
\]
By \eqref{eq:Nq-periodic},
$N_q(a;h)/d=(\ph(q)/q)h_d+O_q(1)$, proving
\eqref{eq:periodic-density-delta-q}.  For a fixed product $m=r e_2$, the
correction is independent of its Boolean partition and therefore costs at
most a constant multiple of
\[
\sum_m \omega_q(m)\frac{2^{\Om(m)}}m<\infty.
\]
On decorated rows the remaining outer factor is bounded by
\[
\sum_{d\mid \operatorname{rad}_{\ge5}(h),\ (d,q)=1}v_{q,h}(d)
=\prod_{\substack{p\mid h\\p\ge5,\ p\nmid q}}
\left(1+\frac1{p-2}\right)
\ll_q\log\log h.
\]
This gives the stated correction ledger.
\end{proof}

For squarefree $m$ put
\begin{equation}\label{eq:Dcal-row}
\boxed{
\Dcal_{m,q}^{(\lambda)}(H)
=2^{-\Om(m)}\sum_{e_1e_2=m}D_\lambda(e_1,e_2;H).
}
\end{equation}
There are only $O_q(1)$ low rows, and all estimates below are uniform over
that finite family, so the superscript $(\lambda)$ may be suppressed after
this declaration.  Deleting the improper endpoint partitions contributes a
separate endpoint correction and is not part of the definition
\eqref{eq:Dcal-row}.  The high-prime weight used with this object is always
$\omega_q$ from \eqref{eq:omegaq}.

\subsection{Boolean fixed-product fibers}

Let $m=e_1e_2$ be squarefree, coprime to $\Lambda_q$, and let $u\bmod m$ be the CRT idempotent selecting the ordered partition $(e_1,e_2)$.  For a fixed low residue $c\bmod\Lambda_q$, the combined progression root is the unique class $t_{u,c}\bmod \Lambda_qm$ satisfying
\begin{equation}\label{eq:crt-root}
t_{u,c}\equiv Hu\pmod m,
\qquad
t_{u,c}\equiv c\pmod{\Lambda_q},
\end{equation}
where $H$ denotes the current reduced kernel height.  Write $\bar\Lambda$ for the inverse of $\Lambda_q\pmod m$ and $\bar m$ for the inverse of $m\pmod{\Lambda_q}$, and define
\[
W(t;m)=\frac1{2^{\Om(m)}}
\sum_{u^2\equiv u\ (m)}\e\!\left(\frac{tu}{m}\right).
\]
CRT gives the exact factorization
\begin{equation}\label{eq:fixed-q-fourier-factorization}
\boxed{
\frac1{2^{\Om(m)}}
\sum_{u^2\equiv u\ (m)}
\e\!\left(\frac{jt_{u,c}}{\Lambda_qm}\right)
=\e\!\left(\frac{jc\bar m}{\Lambda_q}\right)
W(jH\bar\Lambda;m).
}
\end{equation}
Hence
\begin{equation}\label{eq:absolute-fourier}
\left|
\frac1{2^{\Om(m)}}
\sum_u\e\!\left(\frac{jt_{u,c}}{\Lambda_qm}\right)
\right|
=|W(jH\bar\Lambda;m)|.
\end{equation}
The low residue contributes only a unit phase and multiplication by $\bar\Lambda$ permutes the high unit classes.

Passing from idempotents to sign roots $x=2u-1$, define
\[
V(t;m)=\frac1{2^{\Om(m)}}
\sum_{x^2\equiv1\ (m)}\e\!\left(\frac{tx}{m}\right).
\]
The complete-unit second moment from the $q=1$ proof therefore survives unchanged:
\begin{equation}\label{eq:second-moment}
\frac1{\ph(s)}
\sum_{x\in(\Zmod/s\Zmod)^\times}|V(Hjx;s)|^2
=\prod_{\substack{p\mid s\\p\nmid j}}
\frac{p-2}{2(p-1)}.
\end{equation}
This is the key fixed-period observation: fixed $q$ changes the phases, but not the absolute Boolean Fourier coefficients seen by the small box.

\section{Fixed-modulus analytic estimates}\label{sec:analytic-port}

The exact divisor expansion has now reduced the theorem to three analytic estimates.  The small box must retain the Boolean gain that makes its logarithmic exponent strictly less than one; the exterior must preserve the weighted-completion and reciprocal savings after the fixed conductor is adjoined; and the one-dimensional rows must have matching logarithmic main terms.  The complete $q=1$ arguments are proved in the companion manuscript \cite{DorseyPointwise2026}.  Here we state the precise fixed-$q$ changes needed to reuse those arguments.  Throughout, $q$ is fixed, so implicit constants may depend on it.

\subsection{The Boolean small box}

The $q=1$ proof splits a squarefree fixed product $m=sr$ into a smooth part and a rough part, applies Shiu's short-interval Brun--Titchmarsh theorem for nonnegative multiplicative functions \cite[Chapter~3, Theorem~1]{ShiuThesis1980} to the rough factor in a reduced class modulo $s$, and combines \eqref{eq:second-moment} with weighted Cauchy.  In the range used here, Shiu supplies an upper bound for a nonnegative multiplicative function on a dyadic interval in one reduced class modulo $s$, with $s<y^{3/4}$ and interval length $y>x^{1/4}$; the moving rough functions have uniformly bounded prime values and vanish on higher prime powers.  Those are the only features of the theorem used below.

The fixed period $\Lambda_q$ does not enlarge the moving Shiu modulus.  To see this explicitly, expand the bounded periodic weight $\lambda$ by the finitely many residue classes $c\pmod{\Lambda_q}$.  For a dyadic product block $m\asymp M$, write uniquely $m=sr$, with every prime of $s$ at most
\[
Y=\exp((\log M)^{0.99})
\]
and every prime of $r$ larger than $Y$.  On the good range $s\le M^{0.01}$, \eqref{eq:fixed-q-fourier-factorization} and multiplicativity of the Boolean transform give
\begin{equation}\label{eq:fixed-q-shiu-factor}
|W(jH\bar\Lambda;sr)|^2
=|W(jH\bar\Lambda\bar r;s)|^2
 |W(jH\bar\Lambda\bar s;r)|^2
\le |W(jH\bar\Lambda\bar r;s)|^2.
\end{equation}
Grouping the rough factor by its reduced residue class $r\pmod s$ therefore produces exactly the same progression problem modulo $s$ as in the companion Shiu calculation: multiplication by the fixed unit $\bar\Lambda$ merely permutes those classes.  The residue $c\pmod{\Lambda_q}$ contributes only the unit phase in \eqref{eq:fixed-q-fourier-factorization}, and there are $O_q(1)$ such residues.  Thus the rough-factor estimate, the complete-unit second moment \eqref{eq:second-moment}, and the weighted-Cauchy exponent are unchanged.  Deleting the finitely many primes dividing $q$ only removes nonnegative terms from the Shiu majorants.

The only additional counting error comes from replacing the exact weighted root count in \eqref{eq:Dlambda} by its complete-period mean.  For each fixed product and low residue this endpoint error is $O_q(1)$; after the Boolean weights are summed it contributes
\[
O_q(1)+O_q(\log Z),
\]
as recorded in \eqref{eq:periodic-defect-cost}.  This is $o(\log h)$ at the chosen scale.  We may therefore import the companion small-box exponent without any hidden moving-modulus cost.

With the same parameter choices
\[
\sigma=\frac23+0.01,
\qquad
\log Z=4(\log h)^\sigma(1+o(1)),
\]
the fixed-$q$ small-box bound is, for a fixed exponent $K_q$,
\begin{equation}\label{eq:fixed-q-smallbox}
\sum_{\substack{1<m\le Z\\(m,HqR)=1}}
\omega_q(m)2^{\Om(m)}|\Dcal_{m,q}(H)|
\ll_q
(\log\log Z)^{K_q}(\log Z)^{2\sqrt{0.505}}
+O_q(\log Z),
\end{equation}
where $\Dcal_{m,q}$ is the corresponding centered Boolean discrepancy.  Since
\begin{equation}\label{eq:exponent}
\sigma\,2\sqrt{0.505}
=0.9617240306403449\ldots<1,
\end{equation}
the centered fixed-$q$ small box is $o_q(\log h)$ pointwise.

The exact use of $N_q(a;h)$ instead of its mean is the correction isolated in
\cref{lem:fixed-q-row-reduction}; together with the endpoint deletion it costs only
\begin{equation}\label{eq:periodic-defect-cost}
O_q(1)+O_q(\log Z)=o(\log h),
\end{equation}
because the fixed-product endpoint defect is $O_q(1)$ and
\[
\sum_m\omega_q(m)\frac{2^{\Om(m)}}m<\infty.
\]

\subsection{Exterior completion and reciprocity}

For each low class $c\bmod\Lambda_q$, the extra factor
\[
\e_{\Lambda_q}(jc\bar m)
\]
is a fixed-conductor periodic function of the high modulus.  On the unit classes modulo $\Lambda_q$ it decomposes into finitely many fixed character/additive-frequency components.

The weighted-completion input uses a unit-preserving squarefree-removal split.
The divisor $d$ comes from the companion coefficient $\gamma(d)$, so $d$ is
squarefree and $(d,6)=1$; the active support further gives $(d,qR)=1$.  Under
precisely these conditions,
\begin{equation}\label{eq:fixed-q-squarefree-removal}
\boxed{
 \mu^2(dn)\1_{(dn,6R)=1}\1_{(dn,q)=1}
 =\1_{(n,6q)=1}\mu^2(n)\1_{(n,dR)=1}
}\qquad
\bigl[d\ \text{squarefree},\ (d,6qR)=1\bigr].
\end{equation}
Only the high-prime factor $dR$ is expanded by inclusion--exclusion:
\[
 \1_{(n,6q)=1}\mu^2(n)\1_{(n,dR)=1}
 =\1_{(n,6q)=1}
 \left(\sum_{\substack{c\mid dR\\c\mid n}}\mu(c)\right)
 \left(\sum_{\substack{f^2\mid n\\(f,dR)=1}}\mu(f)\right).
\]
After writing $n=cf^2\ell$, every nonzero channel therefore has
\[
 c\mid dR,\qquad (f,dR)=1,\qquad
 (f,6q)=1,\qquad(\ell,6q)=1.
\]
In particular $d,c,f,\ell$ are units modulo $\Lambda_q$ wherever an inverse
modulo $\Lambda_q$ is evaluated.  If a moving reciprocal modulus is present,
its unit restriction is retained as a separate outside indicator through the
same split; it likewise forces every factor that is inverted modulo that moving
modulus to be a unit.  No inclusion--exclusion channel containing a factor $2$
or $3$ is ever asked to carry a $\Lambda_q$-reciprocal phase.

The next statement makes the fixed-conductor completion step explicit without
suppressing the zero Fourier mode of the reciprocal phase.  Let
$\Lambda=\Lambda_q$, let $r$ be a squarefree high modulus with
$(r,\Lambda)=1$, and let $I\subset(U,2U]$ be an interval.  For fixed integers
$A,B$ set
\begin{equation}\label{eq:fixed-conductor-sigma}
\Sigma_{r,\Lambda}^{B}(I;A,R)
=\sum_{\substack{\ell\in I\\(\ell,rR)=1}}
 \omega_q(\ell)\,
 \e_r(A\bar\ell)\,
 \e_{\Lambda}(B\bar\ell).
\end{equation}
The support of $\omega_q$ already forces $(\ell,\Lambda)=1$.

\begin{lemma}[Centered fixed-conductor weighted completion]\label{lem:fixed-conductor-completion}
Put
\[
Q=r\Lambda,
\qquad C=A\Lambda+Br,
\qquad \rho=\frac{Q}{(C,Q)},
\qquad
\mathfrak m_Q(C)=\frac{\mu(\rho)}{\ph(\rho)},
\]
and
\[
M_{r,R}(I)=\sum_{\substack{\ell\in I\\(\ell,rR)=1}}\omega_q(\ell).
\]
Then
\begin{equation}\label{eq:combined-reciprocal-phase}
\e_r(A\bar\ell)\e_{\Lambda}(B\bar\ell)
=\e_Q(C\bar\ell)
\end{equation}
for every supported $\ell$, and
\begin{equation}\label{eq:fixed-conductor-completion}
\boxed{
\Sigma_{r,\Lambda}^{B}(I;A,R)
=\mathfrak m_Q(C)M_{r,R}(I)
+O_q\!\left(
\mathcal E(R)2^{\Om(r)/2}\rho^{1/4}U^{-1/2}(\log 2r)^{1/2}
\right),
}
\end{equation}
where
\[
\mathcal E(R)=\prod_{p\mid R}(1+p^{-1/2}),
\qquad M_{r,R}(I)\ll_q1.
\]
The implied constant is uniform in the endpoints of $I$.  Since $Q$ is
squarefree, the displayed main coefficient is the exact normalized Ramanujan
mean of the reciprocal phase.  Full resonance is exactly
\begin{equation}\label{eq:fixed-conductor-resonance}
\rho=1
\Longleftrightarrow
r\mid A\ \text{and}\ \Lambda\mid B.
\end{equation}
In particular, $\rho>1$ does not by itself imply decay: the Ramanujan-mean
term in \eqref{eq:fixed-conductor-completion} must first be removed or summed.
\end{lemma}

\begin{proof}
Because $(r,\Lambda)=1$, CRT gives
\eqref{eq:combined-reciprocal-phase}.  The phase on the unit group modulo $Q$
has complete mean
\[
\frac1{\ph(Q)}\sum_{x\bmod Q}^{*}\e_Q(C\bar x)
=\frac{c_Q(C)}{\ph(Q)}
=\frac{\mu(\rho)}{\ph(\rho)}.
\]
Subtract this constant before completing an interval.  The centered periodic
function has zero sum on every complete $Q$-period, so only one incomplete
period remains.  Inclusion--exclusion over the resonant prime factors and
completion of that remainder on the reduced modulus $\rho$ give the same
squarefree Kloosterman bound as in the companion argument, now with the zero
Fourier mode absent.  Apply the unit-preserving split
\eqref{eq:fixed-q-squarefree-removal}: only divisors of the high-prime factor
$dR$ enter the $c$-sum, while $\1_{(n,6q)=1}$ stays outside the split.  Thus the
inner $\ell$-sum is always over units modulo $\Lambda$, and $c$ and $f$ are
units modulo $\Lambda$ as well.  Any moving reciprocal unit condition is kept
outside in the same way.  The centered incomplete Kloosterman estimate and the
finite $f$-balance are therefore exactly the companion estimates on each legal
unit channel, with at most an $O_q(1)$ fixed-period factor; this gives the error
term in \eqref{eq:fixed-conductor-completion}.  The untwisted dyadic mass
satisfies $M_{r,R}(I)\ll_q1$ by the elementary logarithmic mean for $\omega_q$,
and restrictions only delete nonnegative terms.  Finally, reducing
$C=A\Lambda+Br$ modulo $r$ and modulo $\Lambda$ proves
\eqref{eq:fixed-conductor-resonance}.
\end{proof}

Thus this fixed-$q$ port uses the companion squarefree-removal and finite
$f$-split architecture only after the reciprocal phase has been centered; it
does not invoke an arbitrary-interval estimate for the uncentered phase.

We use here exactly the companion exterior parameters
\[
L=\log h,
\qquad \sigma=\frac23+0.01,
\qquad K=L^{10},
\qquad W=\exp(L^\sigma),
\]
\[
T_{\rm comp}=\exp(L^{0.60}),
\qquad X=T_{\rm comp}^{3}W^2.
\]
We now record the mean-term ledger that is used in the exterior.  Write
\[
r_A=\frac r{(A,r)},
\qquad
\rho_\Lambda=\frac{\Lambda}{(B,\Lambda)}.
\]
Then $\rho=r_A\rho_\Lambda$ and hence
\begin{equation}\label{eq:fixed-conductor-mean-bound}
|\mathfrak m_Q(C)|\ll_q\frac1{\ph(r_A)}.
\end{equation}
In every actual nonzero Vaaler completion channel, the unit scalings introduced
by squarefree removal do not change the gcd with the moving modulus, and the
moving numerator is $jH$ times such a unit.  Since active moving moduli are
coprime to $H$ and $1\le |j|\le K$, one has
\[
r_A=\frac r{(r,j)}.
\]
For $Y\ge K$, Rankin's trick at exponent $1/2$ gives, uniformly in
$1\le|j|\le K$,
\begin{align}
\sum_{r>Y}\frac{\omega_q(r)}{\ph(r/(r,j))}
&\ll_q
Y^{-1/2}
\sum_{d\mid\operatorname{rad}(j)}\omega_q(d)d^{1/2}\notag\\
&\ll_q
Y^{-1/2}
\prod_{\substack{p\mid j\\p\ge5,\ p\nmid q}}
\left(1+\frac{\sqrt p}{p-3}\right)
=Y^{-1/2}K^{o(1)}.
\label{eq:ramanujan-mean-tail}
\end{align}
Indeed, after $r=ds$ with $d=(r,j)$, the remaining tail is bounded by
\[
\sum_{s>Y/d}\frac{\omega_q(s)}{\ph(s)}
\le
\left(\frac dY\right)^{1/2}
\prod_{\substack{p\ge5\\p\nmid q}}
\left(1+\frac{\sqrt p}{(p-3)(p-1)}\right).
\]
The Euler product is convergent.  Since the other high-modulus mass up to
$h^2$ is $O_q(L)$, the main Vaaler frequencies have total $\ell^1$ mass
$O(\log K)$, and the fixed-period expansion has $O_q(1)$ terms, the aggregate
Ramanujan means from all moving moduli $r>K$ contribute
\begin{equation}\label{eq:ramanujan-mean-global}
O_q\!\left(L\log K\,K^{-1/2+o(1)}\right)=o(1).
\end{equation}
For direct completion the moving modulus is in fact $>X$, so
\eqref{eq:ramanujan-mean-tail} is stronger still.

No modewise completion is used for a moving modulus at most $K$.  In the
direct orientation we return, before expanding the fixed-period factor, to the
underlying sawtooth $\psi(x)=x-\lfloor x\rfloor-\tfrac12$ on the combined modulus $Q=r\Lambda$.  For every
nonconstant sawtooth channel its numerator $C_0$ is nonzero modulo $Q$, and
inversion permutes the units; pairing $x$ with $-x$ therefore gives the exact
complete-period identity
\begin{equation}\label{eq:fixed-conductor-sawtooth-zero}
\sum_{x\bmod Q}^{*}\psi(C_0\bar x/Q)=0.
\end{equation}
The constant channel is the already separated deterministic/full-resonance
piece.  The finite $f$-split applied to the remaining incomplete $Q$-period
then gives, with only a fixed-$q$ change in the constant,
\begin{equation}\label{eq:fixed-conductor-small-modulus}
\ll_q
\mathcal E(R)(\log 2K)^{C_1}(K/U)^{1/2},
\end{equation}
for a fixed absolute $C_1$.  This is super-logarithmically small in the exterior range $U>X$.  Thus the
nonzero Fourier modes are never charged individually for their nonzero
Ramanujan means when the moving modulus is small.

In the reciprocal orientation the companion geometric carve-out already
removes the whole region with moving modulus $u\le K$ before slow/moderate
completion.  We retain that whole carve-out for every fixed-period frequency,
not only for the principal one.  Its band and tail remain respectively
\[
O_q(L^\sigma\log L)=o(L),
\qquad O_q(1),
\]
because the fixed-conductor factors are bounded.  Outside the carve-out the
centered error in \eqref{eq:fixed-conductor-completion} passes through the same
weighted-completion and reciprocal estimates as before; the fixed factor
changes only $O_q(1)$ constants.  The fast reciprocal channels are unchanged:
finite Fourier expansion of the fixed-period decoration adds only linear prime
phases, and Karatsuba's derivative tests used here \cite[Theorem~1]{Karatsuba1971} begin at order at least two.
The Rankin tail and inherited forbidden-prime surcharge likewise change only
by fixed-$q$ constants.

Combining the centered completion error, the exact small-modulus sawtooth
cancellation, the retained reciprocal carve-out, and the negligible global
Ramanujan-mean ledger \eqref{eq:ramanujan-mean-global}, the complete hyperbolic
exterior, including the Vaaler remainder frequencies and the slow
moving-resonance contribution, is
\begin{equation}\label{eq:exterior-q}
\boxed{o_q(\log h)}
\end{equation}
pointwise.

\subsection{Decorated rows}

These rows carry a deterministic divisor of the height.  For a height-divisor
row write $e_1=dr$ with $d\mid h$ and $(d,r)=1$.  All active high primes are
coprime to $q$, so $(d,q)=1$.  Under the exact change of variables $y=dz$,
\begin{equation}\label{eq:decorated-residue}
(a+y,q)=1
\Longleftrightarrow
(d^{-1}a+z,q)=1.
\end{equation}
The low $2,3$ ratio is invariant when both the height and the free coordinate
are multiplied by a high unit modulo $6$.  Therefore a decorated row makes only
the finite residue permutation
\begin{equation}\label{eq:decorated-map}
(a,h)\longmapsto(d^{-1}a,h/d)
\end{equation}
inside the fixed family of endpoint classes.

The exterior estimate at this reduced height is run with the centered
completion just proved, not with an uncentered companion estimate.  The actual
row height changes only the factor $H$ in the moving numerator.  Active moving
moduli remain coprime to that height, so in every Vaaler channel
$r_A=r/(r,j)$ exactly as in \eqref{eq:ramanujan-mean-tail}.  Hence the direct
orientation uses \eqref{eq:fixed-conductor-completion} together with the mean
ledger \eqref{eq:ramanujan-mean-global}; the reciprocal orientation retains the
whole moving-modulus-$\le K$ carve-out and uses the same centered estimate and
mean ledger outside it.  Both statements are uniform for actual row heights at
most the ambient $h$.  The outer decoration norm is $O(\log\log h)$, while
\eqref{eq:ramanujan-mean-global} is $O_q(L\log K\,K^{-1/2+o(1)})=o_q(1)$ and
the centered nonzero-frequency saving is arbitrarily logarithmic.  Therefore
summing all decorated rows still gives $o_q(\log h)$.

\subsection{One-dimensional boundary terms}

The two boundary families are matched before they are estimated.  Exact reflection turns the right boundary into a left boundary with residue $-b$; at divisor-family level, the same reflection leaves only the subfamily supported entirely on active primes dividing $h$, which is smaller by an $O_q(\log\log h)$ estimate.  Define the left free one-anchor sum
\begin{equation}\label{eq:Sqa}
S_{q,a}(h)=\sum_{t\in I_q(a;h)}(\SSer_q(\{0,t\})-1).
\end{equation}
Reflection is exact.  If $u=h-t$, then
\[
(a+t,q)=1
\Longleftrightarrow
(b-u,q)=1
\Longleftrightarrow
(-b+u,q)=1,
\]
and $\SSer_q(\{h,t\})=\SSer_q(\{0,u\})$.  Hence
\begin{equation}\label{eq:reflection-q}
\boxed{
\sum_{t\in I_q(a;h)}(\SSer_q(\{h,t\})-1)
=S_{q,-b}(h).
}
\end{equation}

The divisor families require the corresponding row-level reflection, not only
this free-boundary identity.  For a unit residue $r\pmod q$, write
$\lambda_{q,r,h}$, $N_q(r;h)$, and $\Delta_{q,r}$ for
\eqref{eq:lambdaq}, \eqref{eq:IqNq}, and \eqref{eq:Deltaqea} with
$a$ replaced by $r$, and put
\[
R_{q,h}:=\prod_{\substack{p\mid h\\p\ge5,\ p\nmid q}}p.
\]
Define the residue-sensitive pure-one family
\begin{equation}\label{eq:P1-residue-q}
P_{1,q}^{(r)}(h)
:=C_q^*(h)\sum_{f>1}v_{q,h}(f)\Delta_{q,r}(h;f,1),
\end{equation}
and its pure-height-prime support subfamily
\begin{equation}\label{eq:P1-support-q}
P_{1,q,h}^{\mathrm{supp}}(r)
:=C_q^*(h)
\sum_{\substack{f>1\\f\mid R_{q,h}}}
v_{q,h}(f)\Delta_{q,r}(h;f,1).
\end{equation}
Here $P_{2,q}$ denotes the contribution to \eqref{eq:Tq-expansion} with
$e_1=1<e_2$, and $D_{1,q}$ denotes the contribution with
$e_1=d>1$, $d\mid R_{q,h}$, $e_2>1$, and $(e_2,h)=1$.

\begin{lemma}[Fixed-$q$ divisor-family reflection]\label{lem:reflection-family-q}
Let $d\mid R_{q,h}$ and let $e$ lie on the support of $\omega_q$ with
$(e,h)=1$.  Then
\begin{equation}\label{eq:reflection-row-q}
\boxed{\Delta_{q,a}(h;d,e)=\Delta_{q,-b}(h;de,1),}
\end{equation}
and
\begin{equation}\label{eq:reflection-weight-q}
\boxed{v_{q,h}(d)\omega_q(e)=v_{q,h}(de).}
\end{equation}
Consequently the one-dimensional divisor families satisfy the exact identity
\begin{equation}\label{eq:reflection-assembly-q}
\boxed{
P_{2,q}+D_{1,q}
=P_{1,q}^{(-b)}-P_{1,q,h}^{\mathrm{supp}}(-b),
}
\end{equation}
with
\begin{equation}\label{eq:support-bound-q}
P_{1,q,h}^{\mathrm{supp}}(r)=O_q(\log\log h)
\end{equation}
uniformly over the finitely many compatible unit residues $r\pmod q$.
\end{lemma}

\begin{proof}
Put $z=h-y$.  Reflection is an affine bijection modulo every active low
prime, so
$L_p(\{0,y,h\})=L_p(\{0,z,h\})$ for every active
$p\in\{2,3\}$.
Compatibility $h\equiv b-a\pmod q$ gives
\[
(a+y,q)=1
\Longleftrightarrow (b-z,q)=1
\Longleftrightarrow (-b+z,q)=1.
\]
Hence
$\lambda_{q,a,h}(y)=\lambda_{q,-b,h}(z)$.
The same reflection bijects $I_q(a;h)$ with $I_q(-b;h)$, so
$N_q(a;h)=N_q(-b;h)$ exactly.  If $d\mid h$, then
$d\mid y$ if and only if $d\mid z$; the condition $e\mid h-y$
is exactly $e\mid z$.  Since $(d,e)=1$, these are together equivalent
to $de\mid z$.  Subtracting the identical exact density denominator
therefore proves \eqref{eq:reflection-row-q}.  On the coefficient side,
all primes of $d$ divide $h$ while all primes of $e$ do not, so the local
definitions \eqref{eq:omegaq} and \eqref{eq:vqh} give
\eqref{eq:reflection-weight-q} prime by prime.

The family $P_{2,q}$ is the case $d=1$, $e>1$; $D_{1,q}$ is the case
$d>1$, $e>1$.  After reflection these are exactly the pure-one indices
$f=de>1$ having at least one active high prime not dividing $h$.  The
indices missing from $P_{1,q}^{(-b)}$ are therefore precisely the
$f>1$ supported on $R_{q,h}$, which proves
\eqref{eq:reflection-assembly-q}.

It remains to estimate this support family.  If $f\mid R_{q,h}$ and $H=h/f$,
then $f$ is a unit modulo $\Lambda_q$.  Along the progression $y=fx$,
the low/admissibility weight is a bounded $\Lambda_q$-periodic sequence
whose complete-period mean is $\ph(q)/q$: multiplication by $f$ merely
permutes the residue classes modulo $q$, $2$, and $3$, while the active
low-prime local tables have mean one.  Thus
\[
N_{q,r}(h;f,1)=\frac{\ph(q)}q H+O_q(1).
\]
By the exact periodic count \eqref{eq:Nq-periodic},
$N_q(r;h)/f=(\ph(q)/q)H+O_q(1)$, uniformly in $f$, and hence
\[
\Delta_{q,r}(h;f,1)=O_q(1)
\qquad(f\mid R_{q,h}).
\]
The product $C_q^*(h)$ is bounded for fixed $q$, uniformly in $h$, and on
this support
$v_{q,h}(f)=\prod_{p\mid f}(p-2)^{-1}$.  Therefore
\[
|P_{1,q,h}^{\mathrm{supp}}(r)|
\ll_q
\sum_{f\mid R_{q,h}}v_{q,h}(f)
=\prod_{p\mid R_{q,h}}\left(1+\frac1{p-2}\right)
\ll_q\log\log h,
\]
by the standard Mertens bound, proving \eqref{eq:support-bound-q}.
\end{proof}

For a one-dimensional row let $H$ again be the actual parity-reduced kernel
height and set
\begin{equation}\label{eq:X0-contract}
\boxed{X_0=H-1.}
\end{equation}
The symbol $X$ below is only the current upper limit in an $e$-sum,
$1\le X\le X_0$; it is not a second height.  In the P\'etermann split put
\[
 z=X_0/\Lambda_q,
 \qquad
 Y=z\exp\!\left(- (\log z)^b\right).
\]
The coefficient $v(e)$ in any row is explicitly one of $\omega_q(e)$,
$\eta_q(e)$, or $v_{q,h}(e)$ from
\eqref{eq:omegaq}, \eqref{eq:etaq}, and \eqref{eq:vqh}.

After splitting a one-dimensional root count by $c\bmod\Lambda_q$ and then
by a unit class modulo $\Lambda_q$, character orthogonality reduces every
nonprincipal periodic contribution to a finite linear combination of sums
\begin{equation}\label{eq:petermann-periodic}
\sum_{e\le X}\chi(e)v(e)
\psi\!\left(\frac{X_0}{\Lambda_q e}+\frac r{\Lambda_q}\right),
\end{equation}
where $\chi$ is a fixed character modulo $\Lambda_q$, $r$ is fixed, and
\[\psi(x)=x-\lfloor x\rfloor-\tfrac12.\]

\begin{lemma}[Fixed-conductor character/rational-shift P\'etermann estimate]\label{lem:fixed-conductor-petermann}
Fix $\Lambda\ge1$.  Let $v$ be one of the squarefree nonnegative base
coefficients occurring in the one-dimensional argument after deleting finitely
many primes, and suppose the ordinary $\alpha=0$ P\'etermann hypotheses hold.
For a fixed Dirichlet character $\chi\pmod\Lambda$, fixed integer $r$,
$0<b<1$, and a real reciprocal numerator $T\ge2\Lambda$, put
\[
 Y_T=\frac{T}{\Lambda}
 \exp\!\left(-\bigl(\log(T/\Lambda)\bigr)^b\right).
\]
Then
\[
 \sum_{e\le Y_T}\chi(e)v(e)
 \psi\!\left(\frac{T}{\Lambda e}+\frac r\Lambda\right)
 \ll_{\Lambda,b}
 (\log T)^{2/3}(\log\log(3T))^{4/3}+\log\log(3T).
\]
For $b=3/4$ the same endpoint splitting as in the ordinary P\'etermann argument
gives, uniformly for the partial sums used below,
\begin{equation}\label{eq:petermann-real-numerator}
 \sum_{e\le X}\chi(e)v(e)
 \psi\!\left(\frac{T}{\Lambda e}+\frac r\Lambda\right)
 \ll_{\Lambda}
 (\log(3T))^{3/4}
 +(\log(3T))^{2/3}(\log\log(3T))^{4/3}
 +\log\log(3T)+1
\end{equation}
for $1\le X\le T$.  If $T<2\Lambda$, the corresponding sum has
$O_\Lambda(1)$ terms and is $O_\Lambda(1)$.
\end{lemma}

\begin{proof}
On the supported unit integers, expand the fixed character by finite additive
Fourier inversion,
\[
 \chi(n)=\sum_{a\bmod\Lambda}\widehat\chi(a)\e(an/\Lambda),
 \qquad \sum_a|\widehat\chi(a)|=O_\Lambda(1).
\]
Thus each prime phase in the P\'etermann exponential-sum stage is changed from
$F(p)$ to $F(p)+ap/\Lambda$.  Every derivative test used there starts at order
at least two, so all such derivative hypotheses are unchanged.  The monotone
positive coefficient $pv(p)$ remains outside Abel summation and the character
remains in the already controlled oscillatory prime partial sum; prime powers
$\nu\ge2$ vanish by squarefree support.  Absolute and Cauchy steps are
unchanged.  The reciprocal parameter $T$ is allowed to be real throughout
those derivative estimates.

The rational shift multiplies every nonzero sawtooth Fourier mode by the unit
factor $\e(jr/\Lambda)$ and rescales the reciprocal parameter by the fixed
factor $1/\Lambda$.  With smoothing width $1/T$, a jump crossing for
$e<T/\Lambda$ would require
\[
 T+re<\Lambda e n\le T+re+\frac{\Lambda e}{T}.
\]
The last interval has length less than one.  Since
$\Lambda e n-re=e(\Lambda n-r)$ is an integer, any crossing must satisfy
\[
 e(\Lambda n-r)=\lceil T\rceil.
\]
Thus every crossing index $e$ divides $\lceil T\rceil$ (and obeys one fixed
congruence for the complementary factor).  For either base coefficient used
here,
\[
 \sum_{e\mid\lceil T\rceil}v(e)
 \ll_\Lambda
 \prod_{\substack{p\mid\lceil T\rceil\\p\ge5}}
 \left(1+\frac1{p-3}\right)
 \ll_\Lambda\log\log(3T),
\]
with the same or smaller bound for the $1/(p-2)$ local factor.  This is the
entire crossing bill.  The complementary range $T/\Lambda\le e\le T$ has
$O_\Lambda(1)$ absolute coefficient mass by the ordinary mean formula, and the
usual strip from $Y_T$ to $T/\Lambda$ has mass
$O_\Lambda((\log T)^b)$.  Adding the divisor-supported crossing bill proves the
claims.
\end{proof}

For the moving coefficient, the convolution used below is local and exact.
At each active high prime $p\mid h$ put
\[
 a_p=\frac1{p-3},\qquad b_p=\frac1{p-2},
\]
and define the multiplicative correction $c_{q,h}$ by the Euler factors
\begin{equation}\label{eq:cqh-definition}
\frac{1+b_pz}{1+a_pz}
=1+\sum_{\nu\ge1}c_{q,h}(p^\nu)z^\nu,
\qquad
c_{q,h}(p^\nu)=(b_p-a_p)(-a_p)^{\nu-1},
\end{equation}
at primes $p\ge5$ with $p\nmid q$ and $p\mid h$, and by the factor $1$
at every other prime.  Comparing local generating functions gives exactly
\[
 v_{q,h}=\omega_q*c_{q,h}.
\]
Moreover
\[
1+\sum_{\nu\ge1}|c_{q,h}(p^\nu)|
=1+\frac1{(p-2)(p-4)},
\qquad
\sum_{\nu\ge1}|c_{q,h}(p^\nu)|\,\nu\log p
\ll\frac{\log p}{p^2},
\]
so the convergent product over active high primes gives
\begin{equation}\label{eq:fixed-q-convolution-transfer}
 \sup_h\sum_s|c_{q,h}(s)|(1+\log s)<\infty.
\end{equation}
Deleting the fixed primes dividing $q$ can only decrease these local norms.

For $v=v_{q,h}$ expand the convolution inside the partial sum.  Since every
$s$ in the support of $c_{q,h}$ is a unit modulo $\Lambda_q$,
\begin{align*}
&\sum_{e\le X}\chi(e)v_{q,h}(e)
\psi\!\left(\frac{X_0}{\Lambda_q e}+\frac r{\Lambda_q}\right)\\
&\qquad=\sum_{s\le X}c_{q,h}(s)\chi(s)
\sum_{m\le X/s}\chi(m)\omega_q(m)
\psi\!\left(\frac{X_0/s}{\Lambda_q m}+\frac r{\Lambda_q}\right).
\end{align*}
Apply \cref{lem:fixed-conductor-petermann} to the inner base sum with the real
reciprocal numerator $T=X_0/s$.  When $T<2\Lambda_q$ the inner sum is
$O_q(1)$; otherwise \eqref{eq:petermann-real-numerator} applies.  Its right-hand
side is bounded by the same increasing majorant with $T$ replaced by $X_0$.
The uniform $\ell^1$ bound \eqref{eq:fixed-q-convolution-transfer} therefore
transfers that majorant without any integrality assumption on $X_0/s$.
For $\omega_q$ and $\eta_q$ take $s=1$.  Hence, for every coefficient
$v\in\{\omega_q,\eta_q,v_{q,h}\}$,
\begin{equation}\label{eq:petermann-q}
\begin{aligned}
&\sum_{e\le X}\chi(e)v(e)
\psi\!\left(\frac{X_0}{\Lambda_q e}+\frac r{\Lambda_q}\right)\\
&\qquad=O_q\!\left(
(\log X_0)^{3/4}
+(\log X_0)^{2/3}(\log\log X_0)^{4/3}
+\log\log X_0+1
\right)
=o_q(\log X_0).
\end{aligned}
\end{equation}
uniformly for $1\le X\le X_0$.

The principal periodic coefficient can also be checked locally.  For a residue
$r\pmod q$ define the free one-anchor low-period weight
\begin{equation}\label{eq:rho-free-q}
\rho_{q,r}(y)
=\1_{(r+y,q)=1}
\prod_{\substack{p\in\{2,3\}\\p\nmid q}}L_p(\{0,y\}).
\end{equation}

Put $h=2H$, $X_0=H-1$, and $\Lambda=\Lambda_q$.  The exact finite-period
floor decomposition needed for the one-dimensional bracket is as follows.
Define the two parity-reduced low weights
\begin{equation}\label{eq:one-d-reduced-weights-q}
w^{A}_{r,h}(z):=\lambda_{q,r,h}(2z),
\qquad
w^{F}_{r}(z):=\rho_{q,r}(2z),
\end{equation}
viewed as $\Lambda$-periodic functions.  For
$u\in(\mathbb Z/\Lambda\mathbb Z)^\times$ and $c\bmod\Lambda$, let
$\alpha_u(c)\in\{1,\ldots,\Lambda\}$ be the positive representative of
$c u^{-1}\pmod\Lambda$, with the zero class represented by $\Lambda$, and set
\begin{align}
W(w)&:=\sum_{c\bmod\Lambda}w(c),\label{eq:one-d-W}\\
B_w(u)&:=\sum_{c\bmod\Lambda}w(c)
\left(\frac12-\frac{\alpha_u(c)}\Lambda\right),\label{eq:one-d-B}\\
\Psi_w(u;t)&:=\sum_{c\bmod\Lambda}w(c)
\psi\!\left(\frac t\Lambda+1-\frac{\alpha_u(c)}\Lambda\right).
\label{eq:one-d-Psi}
\end{align}
Then, for every $e$ on the support of the high coefficients and
$u_e\equiv e\pmod\Lambda$,
\begin{equation}\label{eq:exact-periodic-floor-count}
\boxed{
\sum_{1\le n\le X_0/e}w(en)
=\frac{W(w)}\Lambda\frac{X_0}{e}
+B_w(u_e)-\Psi_w(u_e;X_0/e).
}
\end{equation}
Indeed the contribution of the residue $c$ is exactly
\[
w(c)\left\lfloor
\frac{X_0}{\Lambda e}+1-\frac{\alpha_{u_e}(c)}\Lambda
\right\rfloor,
\]
and \eqref{eq:exact-periodic-floor-count} is just
$\lfloor x\rfloor=x-\psi(x)-\tfrac12$ summed over $c$.

Let
\[
N^F_{q,r}(h;e):=
\sum_{\substack{0<y<h\\e\mid y}}\rho_{q,r}(y),
\qquad
\Delta^F_{q,r}(h;e):=N^F_{q,r}(h;e)-\frac{N_q(r;h)}e,
\]
and define the residue-$r$ corner and the free high-prime baseline by
\[
C_q^{(r)}:=C_q^*(h)\Delta_{q,r}(h;1,1),
\qquad
C_{f,q}:=\prod_{\substack{p\ge5\\p\nmid q}}c_1(p).
\]
The exact high-prime Euler expansions give
\begin{align}
Q_{1,q}^{(r)}&:=C_q^{(r)}+P_{1,q}^{(r)}
=C_q^*(h)\sum_{e\ge1}v_{q,h}(e)\Delta_{q,r}(h;e,1),
\label{eq:exact-Q1q}\\
S_{q,r}&=C_{f,q}\sum_{e\ge1}\eta_q(e)\Delta^F_{q,r}(h;e).
\label{eq:exact-Sq}
\end{align}
For \eqref{eq:exact-Sq}, the subtraction is exact because primewise
\[
c_1(p)\left(1+\frac1{p(p-2)}\right)=1,
\]
so $C_{f,q}\sum_e\eta_q(e)/e=1$.  Also
$C_q^{(r)}=O_q(1)$ by the same complete-period mean used in
\cref{lem:fixed-q-row-reduction}.

Substituting \eqref{eq:exact-periodic-floor-count} into
\eqref{eq:exact-Q1q}--\eqref{eq:exact-Sq} gives the literal formulas
\begin{align}
Q_{1,q}^{(r)}
={}&C_q^*(h)\sum_{e\le X_0}v_{q,h}(e)
\left\{
\frac{X_0W(w^A_{r,h})/\Lambda-N_q(r;h)}e
+B_{w^A_{r,h}}(u_e)-\Psi_{w^A_{r,h}}(u_e;X_0/e)
\right\}\notag\\
&-C_q^*(h)N_q(r;h)\sum_{e>X_0}\frac{v_{q,h}(e)}e,
\label{eq:exact-Q1-floor-sawtooth}\\
S_{q,r}
={}&C_{f,q}\sum_{e\le X_0}\eta_q(e)
\left\{
\frac{X_0W(w^F_r)/\Lambda-N_q(r;h)}e
+B_{w^F_r}(u_e)-\Psi_{w^F_r}(u_e;X_0/e)
\right\}\notag\\
&-C_{f,q}N_q(r;h)\sum_{e>X_0}\frac{\eta_q(e)}e.
\label{eq:exact-S-floor-sawtooth}
\end{align}
These are identities, not asymptotic replacements.

For completeness, they decompose into exactly the analytic pieces bounded
above.  On the unit group,
\[
\1_{u_e=u}=\frac1{\ph(\Lambda)}
\sum_{\chi\bmod\Lambda}\overline{\chi(u)}\chi(e),
\]
so the $B$-terms are a finite family of fixed-character coefficient sums and
the $\Psi$-terms are a finite family of the sawtooth sums
\eqref{eq:petermann-periodic}.  The nonprincipal fixed-character coefficient
sums are $O_q(1)$; the principal one has
\[
\sum_{e\le X}\chi_0(e)v(e)=F_0(v)\log X+O_q(1)
\]
for $v=\eta_q$ and, uniformly in $h$, for $v=v_{q,h}$ by the exact convolution
\eqref{eq:fixed-q-convolution-transfer}.  The density terms in braces are
$O_q(1)/e$, since
\[
\frac{W(w^A_{r,h})}{\Lambda}
=\frac{W(w^F_r)}{\Lambda}=\frac{2\ph(q)}q,
\qquad
N_q(r;h)=\frac{2\ph(q)}qH+O_q(1),
\]
and the two tails in
\eqref{eq:exact-Q1-floor-sawtooth}--\eqref{eq:exact-S-floor-sawtooth} are
$O_q(1)$ by the corresponding logarithmic mean and partial summation.

It remains only to identify the complete principal boundary coefficient.  Put
\[
\overline B(w):=\frac1{\ph(\Lambda)}
\sum_{u\in(\mathbb Z/\Lambda\mathbb Z)^\times}B_w(u).
\]
For $c\ne0\pmod\Lambda$, the orbit of $cu^{-1}$ is the unit group modulo
$\Lambda/(c,\Lambda)$; pairing $x$ with its negative shows that the average of
$\alpha_u(c)/\Lambda$ is $1/2$.  For $c=0$ it is identically $1$.  Therefore
\begin{equation}\label{eq:Bbar-zero-residue}
\boxed{\overline B(w)=-\frac12w(0).}
\end{equation}
At the zero low residue the anchored and free weights agree exactly:
for every active $p\in\{2,3\}$,
\[
\frac{L_p(\{0,0,h\})}{L_p(\{0,h\})}
=(1-1/p)^{-1}=L_p(\{0,0\}),
\]
and the $q$-admissibility indicator is $1$.  Hence
\begin{equation}\label{eq:Bbar-anchored-free}
\overline B(w^A_{r,h})=\overline B(w^F_r).
\end{equation}
CRT and the low-prime local tables give the exact complete-period means
\begin{equation}\label{eq:principal-low-mean}
\frac1{\Lambda_q}\sum_{c\bmod\Lambda_q}\lambda_{q,r,h}(c)
=\frac1{\Lambda_q}\sum_{c\bmod\Lambda_q}\rho_{q,r}(c)
=\frac{\ph(q)}q.
\end{equation}
Indeed every prime $p\mid q$ contributes the common factor $1-1/p$; at an
active prime $2$ the two local weights have values $(2,0)$ and mean one; at
an active prime $3$ the free values are $(3/2,3/4,3/4)$, while the anchored
values are the same when $3\mid h$ and, when $3\nmid h$, are a permutation
of $(3/2,3/2,0)$.  Each active low-prime mean is therefore one.

For a squarefree coefficient $u$ occurring here, write its logarithmic Euler
residue as
\[
F_0(u):=\prod_p\left(1-\frac1p\right)(1+u(p)),
\]
and recall the free high-prime baseline $C_{f,q}$ defined above.
The active high-prime identities
\[
 c_0(p)\left(1-\frac1p\right)\left(1+\frac1{p-3}\right)=1,
 \qquad
 c_1(p)\left(1-\frac1p\right)\left(1+\frac1{p-2}\right)=1
\]
then give the literal coefficient equality
\begin{equation}\label{eq:principal-high-coefficient}
C_q^*(h)F_0(v_{q,h})
=C_{f,q}F_0(\eta_q)
=\prod_{p\mid6q}\left(1-\frac1p\right)
=\frac{\ph(\Lambda_q)}{\Lambda_q}.
\end{equation}
The exact floor/sawtooth decompositions
\eqref{eq:exact-Q1-floor-sawtooth}--\eqref{eq:exact-S-floor-sawtooth} now leave
no unassigned one-dimensional term.  Their finite sawtooth character families, including the principal character,
are $o_q(\log h)$ by \eqref{eq:petermann-q}; the nonprincipal
constant-character pieces and all density, endpoint, and tail terms are
$O_q(1)$.  The only logarithmic principal terms are
\[
C_q^*(h)F_0(v_{q,h})\overline B(w^A_{r,h})\log X_0
\quad\text{and}\quad
C_{f,q}F_0(\eta_q)\overline B(w^F_r)\log X_0.
\]
They are equal by \eqref{eq:principal-high-coefficient} and
\eqref{eq:Bbar-anchored-free}.  Thus
$Q_{1,q}^{(r)}-S_{q,r}=o_q(\log h)$; since
$C_q^{(r)}=O_q(1)$, the residue-sensitive notation
\eqref{eq:P1-residue-q} gives
\begin{equation}\label{eq:pure-one-uniform-q}
P_{1,q}^{(r)}-S_{q,r}=o_q(\log h)
\end{equation}
for every compatible unit residue $r\pmod q$ occurring in the reflected pair.
Taking $r=a$ proves the first bracket.  Taking $r=-b$ and then applying the
exact family reflection \eqref{eq:reflection-assembly-q} and support bound
\eqref{eq:support-bound-q} proves the second.  Explicitly,
\begin{equation}\label{eq:one-dimensional-q}
P_{1,q}^{(a)}-S_{q,a}=o_q(\log h),
\qquad
P_{2,q}+D_{1,q}-S_{q,-b}
=\bigl[P_{1,q}^{(-b)}-S_{q,-b}\bigr]
-P_{1,q,h}^{\mathrm{supp}}(-b)
=o_q(\log h).
\end{equation}

\section{Proof of the fixed-modulus theorem}\label{sec:proof-main}

Only assembly remains.  The exact divisor expansion \eqref{eq:Tq-expansion} has four geometrically different pieces: the corner, two one-dimensional boundary families, the ordinary two-dimensional rows, and the rows decorated by height primes.  Let $C_q$ be the corner $(e_1,e_2)=(1,1)$; let $P_{1,q}=P_{1,q}^{(a)}$ and $P_{2,q}$ denote respectively the rows $e_2=1<e_1$ and $e_1=1<e_2$; let $G_q$ denote the genuinely two-dimensional rows $e_1,e_2>1$ with $(e_1,h)=1$; and let $D_q$ denote those with $(e_1,h)>1$.  Write $e_1=dr$, where $d=(e_1,h)$.  Split
\[
D_q=D_{1,q}+D_{2D,q}
\]
according as $r=1$ or $r>1$.  Only $D_{1,q}$ is one-dimensional; $D_{2D,q}$ is a decorated two-dimensional family handled by \eqref{eq:decorated-map}.

The exact family assembly is
\begin{equation}\label{eq:final-assembly-q}
\boxed{
A_{1,q}-B_{1,q}
=[P_{1,q}-S_{q,a}]
+[P_{2,q}+D_{1,q}-S_{q,-b}]
+C_q+[G_q+D_{2D,q}].
}
\end{equation}
The corner and the exact periodic-density corrections from
\eqref{eq:fixed-q-row-reduction} contribute
\[
O_q(1)+O_q(\log\log h).
\]
The two one-dimensional brackets are $o_q(\log h)$ by \eqref{eq:one-dimensional-q}.  The centered two-dimensional small box is $o_q(\log h)$ by \eqref{eq:fixed-q-smallbox}--\eqref{eq:exponent}; its hyperbolic fill and the full exterior are $o_q(\log h)$ by \eqref{eq:exterior-q}; and the decorated rows remain $o_q(\log h)$ by \eqref{eq:decorated-map}.  Inserting these estimates into \eqref{eq:final-assembly-q} proves \cref{thm:fixed-q}.

\begin{remark}[Asymptotic character]\label{rem:asymptotic}
The strict Boolean margin is inherited from the $q=1$ proof and is numerically small.  The reciprocal argument also uses a superlogarithmic but very slowly emerging saving.  The theorem is therefore asymptotic; no computationally useful finite-height threshold is asserted.
\end{remark}

\section{Finite-cutoff arithmetic diagonal contraction}\label{sec:collision}

The level-one transfer theorem is already proved.  This section records a separate finite-cutoff identity for the active high-prime coefficient system.  Boolean M\"obius inversion is an incidence-algebra operation; by contrast, averaging a repeated free coordinate after connectedization produces an arithmetic operator on the modulus coefficients.

Fix a finite active high-prime set $\mathcal P$ and an admissible anchor pattern
$A'\subseteq E$.  Free coordinates remain labeled even when two coordinate
values coincide.  For $p\in\mathcal P$ define
\begin{equation}\label{eq:collision-local}
\phi_{p,z}(x)=\frac{\1_{x\not\equiv z\ (p)}}{1-1/p},
\qquad
L_p(\mathbf x)=\frac1p\sum_{z\bmod p}\prod_j\phi_{p,z}(x_j),
\end{equation}
and
\begin{equation}\label{eq:collision-gw}
g_p^{A'}(\mathbf x)=\frac{L_p(A'\cup\mathbf x)}{L_p(A')},
\qquad
w_p^{A'}(\mathbf x)=g_p^{A'}(\mathbf x)-1.
\end{equation}
For squarefree $m\mid\prod_{p\in\mathcal P}p$ set
\begin{equation}\label{eq:collision-u}
u_1^{A'}=1,
\qquad
u_m^{A'}(\mathbf x)=\prod_{p\mid m}w_p^{A'}(\mathbf x).
\end{equation}
Then the truncated relative Euler product is reconstructed exactly by
\begin{equation}\label{eq:collision-euler-reconstruction}
\sum_{m\mid\prod_{p\in\mathcal P}p}u_m^{A'}(\mathbf x)
=\prod_{p\in\mathcal P}g_p^{A'}(\mathbf x).
\end{equation}
For a labeled set $J$ define
\begin{equation}\label{eq:collision-mob}
\operatorname{Mob}_JF
=\sum_{V\subseteq J}(-1)^{|J|-|V|}F(\mathbf x_V).
\end{equation}
For a doubled block $B=\{1,2\}$ define diagonal contraction at modulus $m$ by
\begin{equation}\label{eq:collision-D}
(\mathsf D_BF)_m(\mathbf y)
=\frac1m\sum_{x\bmod m}F_m(x,x,\mathbf y).
\end{equation}
Finally put
\begin{equation}\label{eq:L-operator}
(\Kcal f)_m=\frac1{\ph(m)}\sum_{d\mid m}d f_d,
\qquad
\Lcal=\Kcal-I.
\end{equation}

\begin{proposition}[Finite-cutoff collision intertwiner]\label{prop:collision-intertwiner}
For every squarefree finite-cutoff modulus $m$,
\begin{equation}\label{eq:collision-intertwiner}
\boxed{
\mathsf D_B\operatorname{Mob}_{B\cup U}u_m^{A'}
=\operatorname{Mob}_U(\Lcal u^{A'})_m.
}
\end{equation}
\end{proposition}

\begin{proof}
Write
\[
G_d^{A'}(\mathbf x)=\prod_{p\mid d}g_p^{A'}(\mathbf x).
\]
Then inclusion--exclusion in the active prime set gives
\begin{equation}\label{eq:collision-uG}
u_m^{A'}
=\sum_{d\mid m}(-1)^{\Om(m/d)}G_d^{A'}.
\end{equation}
Under the anchor-conditioned product measure, $G_d^{A'}$ is an expectation
of products of
\[
\Phi_{d,\mathbf z}(x)=\prod_{p\mid d}\phi_{p,z_p}(x).
\]
Boolean connectedization therefore produces a product of
$\Phi_{d,\mathbf z}-1$.  CRT gives the exact diagonal second moment
\begin{equation}\label{eq:collision-second-moment}
\frac1d\sum_{x\bmod d}(\Phi_{d,\mathbf z}(x)-1)^2
=\frac d{\ph(d)}-1,
\end{equation}
independently of the anchor-conditioned $\mathbf z$.  Hence
\begin{equation}\label{eq:collision-G}
\mathsf D_B\operatorname{Mob}_{B\cup U}G_d^{A'}
=\left(\frac d{\ph(d)}-1\right)\operatorname{Mob}_UG_d^{A'}.
\end{equation}
Substituting \eqref{eq:collision-uG} into \eqref{eq:L-operator}, the coefficient
of $G_d^{A'}$ in $(\Lcal u^{A'})_m$ is exactly
\[
(-1)^{\Om(m/d)}\left(\frac d{\ph(d)}-1\right).
\]
Comparison with \eqref{eq:collision-G} proves
\eqref{eq:collision-intertwiner}.
\end{proof}

Define the anchor-combined rank-$r$ transfer kernel by
\begin{equation}\label{eq:collision-Dr}
D_m^{(r)}(\mathbf x)
=\operatorname{Mob}_{[r]}
\sum_{A'\subseteq E}(-1)^{|E|-|A'|}u_m^{A'}(\mathbf x).
\end{equation}
Because $\Lcal$ is independent of the anchor pattern, it commutes with this
anchor combination and gives
\begin{equation}\label{eq:collision-one}
\frac1m\sum_{x\bmod m}D_m^{(r+2)}(x,x,\mathbf y)
=(\Lcal D^{(r)})_m(\mathbf y).
\end{equation}
Iterating disjoint doubled blocks gives
\begin{equation}\label{eq:collision-iterate}
\frac1{m^j}\sum_{x_1,\ldots,x_j\bmod m}
D_m^{(2j+1)}(x_1,x_1,\ldots,x_j,x_j,y)
=(\Lcal^jD^{(1)})_m(y).
\end{equation}

Equations \eqref{eq:collision-intertwiner}--\eqref{eq:collision-iterate} are
finite-cutoff identities for the active high-prime system alone.  They are not
used in the proof of \cref{thm:fixed-q}, and no asymptotic or rank-uniform
statement is attached to them here.

\section{Relation to the completed fixed-gap density}\label{sec:los}

The contraction identity \eqref{eq:F-Rcirc} is the exact bridge to the completed endpoint coordinates used in the fixed-gap consecutive-prime residue formalism.  At level one,
\[
C_{1,q}^{\mathrm{comp}}=\SSer_q(E)A_{1,q},
\]
while the transfer theorem controls the connected defect between this anchored coordinate and the sum of the two one-anchor connected boundaries.  The ambient modulus density is carried separately by \eqref{eq:deltaq}; because $\delta_q\ph(q)/q=1/\log x$, there is no hidden $q/\ph(q)$ renormalization inside the connected transfer kernel.  Nothing in this note supplies a rank-uniform tail estimate or rederives a Poisson survival factor.

\section{Source and priority note}\label{sec:priority}

The mathematical assertions of \cref{prop:endpoint-contraction,thm:fixed-q} are established by the proofs above.  Priority standing is separate; proof standing is not used as evidence of novelty or priority.

We prove the exact contraction dictionary and the fixed-$q$ pointwise connected transfer theorem stated above. A 2026 preprint of Hamed Ebadi, referenced by a public Zenodo verification record \cite{Ebadi2026}, studies the same fixed-two-anchor raw coefficient hierarchy at $q=1$. The manuscript proof was not available to us for inspection at the time of this version, so we make no claim here about priority relative to proof-level lemmas in that preprint beyond the literal distinctions supported by its public theorem statement.

For the same reason, this note makes no novelty claim for fixed anchors, congruence restrictions, Boolean M\"obius inversion, generic connected/cumulant language, or the conditional fixed-gap Poisson factor.  Those issues belong to the prior-art record rather than to the proof of \cref{thm:fixed-q}.

\begin{thebibliography}{9}

\bibitem{DorseyPointwise2026}
C.~Dorsey,
\emph{Boolean CRT Fibers and Pointwise Anchor Contraction},
preprint manuscript, September 2026.

\bibitem{ShiuThesis1980}
P.~M.-K.~Shiu,
\emph{Studies in Multiplicative Number Theory},
Ph.D. thesis, University of Nottingham, May 1980, Chapter~3.

\bibitem{Karatsuba1971}
A.~A.~Karatsuba,
\emph{Estimates for trigonometric sums by Vinogradov's method, and some applications},
Proc. Steklov Inst. Math. 112 (1971), 251--265
(English translation, American Mathematical Society, 1973).

\bibitem{Petermann1998}
Y.-F.~S.~P\'etermann,
\emph{On an estimate of Walfisz and Saltykov for an error term related to the Euler function},
J. Th\'eor. Nombres Bordeaux 10 (1998), no.~1, 203--236.

\bibitem{GoldstonLedoan2011}
D.~A.~Goldston and A.~H.~Ledoan,
\emph{On the differences between consecutive prime numbers, I},
arXiv:1111.3380.

\bibitem{LOS}
R.~J.~Lemke Oliver and K.~Soundararajan,
\emph{Unexpected biases in the distribution of consecutive primes},
preprint manuscript.

\bibitem{Kuperberg2025}
V.~Kuperberg,
\emph{Sums of singular series along arithmetic progressions and with smooth weights},
arXiv:2301.06095v1 (2023).

\bibitem{Ebadi2026}
H.~Ebadi,
2026 preprint referenced by Zenodo verification record 21534407.
Manuscript proof unavailable to the author of this version at the time of writing.

\end{thebibliography}

\end{document}
